\documentclass[11pt,a4paper]{article}
\usepackage[T1]{fontenc}
\usepackage[utf8]{inputenc}
\usepackage{lmodern}
\usepackage[margin=2.5cm]{geometry}
\usepackage[hidelinks]{hyperref}
\usepackage{parskip}

\begin{document}
\pagestyle{empty}

\noindent\textbf{Article type:} Original Research\\
\textbf{Title:} Within-Family Gains vs External Skill in Fractional SEIT Tuberculosis Forecasts\\
\textbf{Authors:} Amaury de Souza and Francisco Federico Garc\'{\i}a Crespi\\
\textbf{Corresponding author:} Francisco Federico Garc\'{\i}a Crespi (\href{mailto:fedeg@umh.es}{fedeg@umh.es})

\hrule
\bigskip

Dear Editor-in-Chief,

Please consider our manuscript, ``Within-Family Gains vs External Skill in Fractional SEIT Tuberculosis Forecasts,'' for publication as an Original Research article in the \emph{Bulletin of Mathematical Biology}.

The manuscript asks a methodological question that is easy to overlook in applications of fractional-order epidemic models: when a fractional-order model fits and forecasts better than its integer-order counterpart, does that advantage survive when the model is compared with simple forecasts that do not use a mechanistic model? We address it for a fractional-order susceptible-exposed-infectious-treated (SEIT) model of tuberculosis, using 264 monthly notifications from Brazil (2001--2022). We reimplemented the model independently, refitted the integer-order version separately under the same calibration conditions, and compared both with persistence, seasonal naive and SARIMA forecasts over 13 expanding-window origins and leads of 1 to 12 months.

Fractional SEIT had lower RMSE and MAE than integer SEIT at every lead (RMSE at lead 1: 792 against 1186 cases per month). Persistence (390) and SARIMA (536) had lower errors than both SEIT variants at every lead, and the fractional model improved on seasonal naive (1091) only at leads 1 to 5. We therefore conclude that, for this series and window, the within-family improvement does not amount to competitive forecasting performance against the strongest simple benchmarks.

We state the limits of that conclusion in the paper. The comparisons are descriptive: there are 13 origins, the errors overlap, the evaluation window follows the disruption of 2020, and no formal test of equal predictive accuracy was run. The fitted fractional order and progression rate lie at their lower search bounds. A sensitivity analysis with wider bounds, which we report as such, moves both estimates and narrows the within-family margin, without changing its sign. A secondary analysis shows that individual transmission, removal and excess-mortality rates are weakly identified while predicted trajectories are similar, and that the reference reproduction number and disease-free equilibrium describe an autonomous reference, not the time-varying forecasting system. The manuscript makes no claim of general fractional-model superiority, of biological memory, or of priority.

We submit to the \emph{Bulletin} because the question sits at the interface of mathematical modelling and its empirical assessment. The journal has published work on fractional-order epidemic models (Angstmann, Henry and McGann, 2016) and on prediction from poorly identified models (Simpson and Maclaren, 2024). Our manuscript contributes an empirical evaluation to that literature; it does not claim a theoretical advance.

\emph{Transparency.} The code, the numerical results and the frozen evidence on which every reported number rests are archived at Zenodo (DOI: \href{https://doi.org/10.5281/zenodo.23298329}{10.5281/zenodo.23298329}), and the repository is public. The monthly series was supplied by the first author; its original extraction settings could not be recovered, and a pre-specified audit of four DATASUS/SINAN reconstructions did not reproduce it exactly (the closest differed in 99 of 264 months). The manuscript states this in Methods, Limitations and Data Availability. The code was generated with the help of artificial intelligence tools and reviewed by the authors, as declared in the manuscript.

The authors declare no competing interests and received no funding. The study used only secondary, aggregated notification data. The manuscript has not been published and is not under consideration elsewhere, and both authors have approved the final version.

Thank you for considering the manuscript.

Sincerely,

\bigskip

Francisco Federico Garc\'{\i}a Crespi (corresponding author)\\
Department of Computer Engineering, Universidad Miguel Hern\'andez de Elche, Elche, Spain\\
\href{mailto:fedeg@umh.es}{fedeg@umh.es} \textperiodcentered{} ORCID 0000-0002-7129-7791

On behalf of both authors (Amaury de Souza, Instituto de F\'{\i}sica, Universidade Federal de Mato Grosso do Sul, Campo Grande, Brazil; ORCID 0000-0001-8168-1482)

\end{document}
